\documentclass[11pt]{article}
\usepackage[a4paper,margin=27mm]{geometry}
\usepackage{amsmath,amssymb,amsthm}
\usepackage[hidelinks]{hyperref}
\newtheorem{theorem}{Theorem}[section]
\newtheorem{lemma}[theorem]{Lemma}
\newtheorem{definition}[theorem]{Definition}
\newtheorem{corollary}[theorem]{Corollary}
\theoremstyle{remark}
\newtheorem{remark}[theorem]{Remark}
\newcommand{\id}{\mathrm{id}}
\newcommand{\Phib}{\overline{\Phi}}
\newcommand{\Cst}{\mathrm{Cst}}
\newcommand{\Acc}{\mathrm{Acc}}
\newcommand{\Out}{\mathrm{Out}}
\newcommand{\rep}{\mathrm{rep}}
\renewcommand{\tt}[1]{\texttt{#1}}
\title{Cut elimination for $\mathrm{ID}_1$ in $\beta$-logic:\\ a constructive reading of Girard's construction}
\author{}
\date{21 November 2026}
\begin{document}
\maketitle

\begin{abstract}
We report on a machine-checked, constructive reconstruction of the core of
Girard's cut elimination for the inductive logic of $\mathrm{ID}_1$ in
$\beta$-logic \cite[\S11.4]{ptlc2}. The input is a homogeneous family of
proofs, presented by a functor on well-orders, and the induction is on an
explicit \emph{accessibility certificate} of this family. For inputs whose
cuts are atomic (of any polarity and stage), or compound with every
$X$-atom reading a stage below $\mathrm{top}_Y$, we obtain at every well-order $Y$ a cut-free
derivation on $kY+E$, where the well-order $E$ is itself an output of the
construction. Embeddings of the parameter are carried to explicit
embeddings of these outputs. The central point is Girard's treatment of
the cut on the completed predicate $\Phib$: the positive premise is
normalised first, and its output determines the larger order at which the
negative premise is evaluated. All results below are checked in the Unit
development (files \tt{BetaIDS4*.seg}, \tt{BetaAcc.seg}); no excluded
middle and no choice are used, except where stated.
\end{abstract}

\section{The problem and the setting}

For a positive operator $\Phi(X,x)$, cut elimination for $\mathrm{ID}_1$
in $\omega$-logic (Pohlers, Buchholz, Sieg) works when the inductive
predicate occurs only positively in the conclusion. Girard's
$\Pi^1_2$-logic, or $\beta$-logic, gives full cut elimination with a
genuine subformula property, at the price of a proof theory built on
dilators \cite[ch.~11]{ptlc2}. Akiyoshi and Mints \cite{am} obtain complete
cut elimination for arbitrary sequents in $\omega$-logic with an extended
$\Omega$-rule, but without a subformula property: the rule for
$\neg\forall X$ admits an arbitrary instance. Our question is how much of
Girard's construction is constructive, and where its strength
(11.6.4: cut elimination implies $\Pi^1_1$-$\mathrm{CA}$) actually enters.

\paragraph{Syntax.} Formulas are built from decidable atoms, $\wedge$,
$\vee$, number quantifiers, and $X$-atoms whose predicate is read from an
environment label: either $\Phib$ (the completed fixed point) or a stage
$I^\mu$ with $\mu$ in a linear order. Rules are the logical rules, the
$x$-rules $(xr)$ (introduce $I^\mu(t)$ from $\Phi(I^{<\mu},t)$, with a
choice $\nu<\mu$ recorded as a point) and $(xl)$ (an $\omega$- and
ordinal-branching rule with a bound), the barred rules for $\Phib$, and
cuts. Following Girard's Definition~11.4.2, a cut is of \emph{kind I}
if its formula has no $X$-atom, and of \emph{kind III} if it is on
$\Phib(t)$ (either polarity in our one-sided calculus). All remaining
cuts are of \emph{kind II}, including cuts on stage atoms and cuts on
compound formulas containing $X$-atoms.

\paragraph{Two standing hypotheses.} Atoms depend only on the numbers of
the labels (\tt{EvN}). The operator satisfies $\mathrm{posX}\,\phi\,0$: the
$X$-atoms of $\phi$ are positive, read the stage position, and never use
it as an argument. In a cut-free derivation, if negative stages in the
end sequent lie below top, this bound propagates to the active premises.
This is what makes Girard's transfer
lemma 11.4.6 go through (\tt{BetaIDCfOps.Inv}, \tt{TrL}).

\paragraph{Input families.} An input is a homogeneous $F$-proof, presented
as a record \tt{IObj}:
\begin{itemize}
\item a functor $k$ on linear orders, with maps $k(f):kX\to kY$ for
embeddings $f$;
\item a decomposition $\sigma_Y : kY\cong(\Cst+Y)+(1+F'Y)$, natural below
top, with $\mathrm{top}_Y=\sigma_Y^{-1}(\mathrm{inr}\,(\mathrm{inl}\,\ast))$ and
$k(f)(\mathrm{top}_X)=\mathrm{top}_Y$;
\item at each $Y$, an end sequent $G_Y$ and a rule function $R_Y$ on
addresses, natural in $Y$, locally correct, with every cut of one of two classes
(\texttt{Obj.KindOK}): \emph{below top}, every $X$-atom of the cut formula
reads a stage $\mu<\mathrm{top}_Y$, never $\Phib$ (\texttt{Obj.StLow};
compound formulas and number quantifiers are allowed, and kind I is the
case without $X$-atoms); or \emph{atomic represented as $\Phib$}, an
$X$-atom reading $\Phib$ or a stage $\mu\ge\mathrm{top}_Y$
(\texttt{Obj.BarR}; kind III is the first case). Together: every atomic
cut, whatever its polarity and stage, and every compound cut below top.
\end{itemize}
No functor laws for $k$ are assumed. The decomposition $\sigma$ replaces
Girard's subfunctor bookkeeping: shifting by an order $D$ gives the family
$k(D+-)$ with constants $\Cst+D$, and $\sigma$ is re-associated.

\paragraph{Accessibility instead of correctness.} Girard proves 11.4.4 by
induction on $\pi$ along the predecessor relation of 10.2.9. A predecessor
of a family is either a plain step (a natural label at the root) or a
shifted step (a variable ordinal premise, the shift by $D$). Constructively
we take the inductive accessibility predicate along all predecessors as
\emph{data}: a well-founded tree whose nodes are the predecessors, with a
root equation (\tt{BetaAcc.AccR}, \tt{AccP}).

\begin{theorem}[\tt{BetaAcc.accCorr}]
An accessible object whose functor is functorial and whose trees are
prefix-closed is correct: its tree at every well-order has a Brouwer
bound. The proof is constructive.
\end{theorem}

The trick is that no splitting of the order is needed: a node $a$ at $X$
is the predecessor with shift $D=X$ and label $k(\mathrm{inl})(a)$, and
evaluated at the empty order it is the subtree at $a$. The converse is
Girard's 10.2.9 and is classical.

\begin{theorem}[\tt{BetaAcc.ClassicalAcc.corrAcc}]
Assuming excluded middle and dependent choice (both at the level of
large types), every correct object has an accessibility certificate.
\end{theorem}

An earlier version also needed choice over the predecessors, to collect
certificates of all children. That hypothesis is gone because certificates
are unique (Theorem~\ref{thm:accprop} below), so the bar argument of
\tt{BetaPred.Cl} can eliminate directly into $\Acc$. We read this as
locating the $\Pi^1_1$-$\mathrm{CA}$ content of Girard's theorem in the
bar theorem ``correct $\Rightarrow$ accessible''. The rest of this note is
about what one gets once a certificate is given.

\section{The pointwise theorem}

\paragraph{Outputs.} Fix an input $o$ and a well-order $Y$. Put
$A=kY$, $\mathrm{top}=\mathrm{top}_Y$. An \emph{output} is a well-order $E$ together with
a cut-free derivation (an inductive object, \tt{Cf}) of $\rep(G_Y)$ in the
context $(A+E,\ \mathrm{inl}\,\mathrm{top})$. Here $\rep$ replaces every stage
label $\ge\mathrm{top}$ by $\Phib$ and keeps the stages below top; this is
Girard's $\widetilde\Gamma$. The extension $E$ lies above top: it collects
the ordinal material created by normalisation.

\begin{theorem}[\texttt{BetaIDS4.Thm.thm}]\label{thm:pointwise}
Let $o$ be an input family satisfying the cut condition above, and let
$ac$ be an accessibility certificate for $o$. For every well-order $Y$
there is a well-order $E$ and a cut-free derivation of $\rep(G_Y)$ on
$(kY+E,\mathrm{inl}\,\mathrm{top}_Y)$. No excluded middle or choice is used.
\end{theorem}

The proof is by recursion on the certificate, with a case split on the
root rule. The rule is read off at the empty order, so that the premise
steps are natural at every order.
\begin{itemize}
\item \emph{Kept rules} (axioms, logical rules, $(xr)$ with choice below
top, $(xl)$, kind I cuts). The premise outputs live on $A+E_a$, one for
each label $a$ of $A$. They are combined over the dependent
lexicographic sum $B=A+\sum_{a}E_a$ (Girard's integral of dilators,
\tt{DSumLO}), transported to $B$ along maps that are initial only below
top (Lemma 11.4.6), and the rule is applied. For a kind I cut, the cut
between cut-free derivations is then eliminated (\texttt{BetaIDCfCut}). The
premises at constant stages are plain steps. The premises at variable
stages $\nu=\sigma^{-1}(\mathrm{inl}(\mathrm{inr}\,y))$ are shifts by
$D=Y$, evaluated at the empty order, and moved back along the
$\sigma$-map $k(Y+0)\to kY+F'(Y+0)$ (\tt{OutBack}).
\item \emph{$(\overline{r})$ above top.} If the choice $\mu$ of an
$(xr)$ rule is $\ge\mathrm{top}$, the premise output proves $\Phi(\Phib,t)$. We add
one point $c$ on top of the extension, rename the marked $\Phib$ into
$I^c$ (Lemma 11.4.5, compound case) and apply $(\overline r)$ with
choice $c$ (\tt{RBar}).
\item \emph{Kind III cuts.} This is the main case, described next.
\end{itemize}

\subsection{The cut on the completed predicate}

Consider a cut on $\Phib(t)$, with positive premise $\pi'$ (proving
$\Gamma,\Phib(t)$) and negative premise $\pi''$ (proving
$\Gamma,\neg\Phib(t)$). In Girard's construction the negative premise is
not evaluated at the same order as the positive one. The idea is to
normalise the positive premise \emph{first}.
\begin{enumerate}
\item The induction hypothesis at $(\pi',Y)$ gives an output $E_1$ and a
cut-free $\widetilde L$ on $A_1=kY+E_1$.
\item Put
\[
 V_1 \;=\; Y+\bigl((1+F'Y)+E_1\bigr).
\]
By re-association, $A_1=(\Cst+Y)+(1+F'Y)+E_1$ is $\Cst+V_1$. So the whole
of $A_1$, including the new material $E_1$ above top, becomes a part of
the \emph{parameter} of the negative premise.
\item The induction hypothesis at $(\pi'',V_1)$ gives an output $E_2$ and
a cut-free $\widetilde R$ on $A_2=kV_1+E_2$.
\item The embedding $g=\mathrm{inl}\circ\sigma_{V_1}^{-1}\circ
\mathrm{inl}\circ\alpha\circ(\sigma_Y+\id):A_1\to A_2$ is initial below
top, and the context of $\widetilde R$ is the $g$-image of that of
$\widetilde L$: below top by naturality of $\sigma$, at and above top both
sides are $\Phib$.
\item Now the stage $\mathrm{top}_{V_1}$ lies above all of $A_1$. The $\Phib$-cut
between $\widetilde L$ (transported along $g$) and $\widetilde R$ becomes
a cut on the stage atom $I^{\mathrm{top}_{V_1}}(t)$ (Lemmas 11.4.5 and 11.4.7),
and that cut is eliminated between cut-free derivations by recursion on
the stage and then the formula (\tt{S2C}, \texttt{BetaIDCfCut}).
\item Finally $A_2=kY+(E_1+((1+F'V_1)+E_2))$ by explicit isomorphisms, so
the result is an output at $Y$ with $E_3=E_1+((1+F'V_1)+E_2)$.
\end{enumerate}
This is the constructive content of Girard's treatment: there is no fixed
order large enough in advance. The order at which the negative premise is
read is \emph{computed} from the normal form of the positive premise. The
recursion is legitimate because both $\pi'$ and $\pi''$ are predecessors
in the accessibility certificate, and the induction hypothesis is stated
for \emph{all} well-orders.

\begin{remark}
Correctness of the outputs is automatic, since cut-free derivations are
inductive objects. The induction hypothesis is used only on input
predecessors, so the outputs need not be accessible themselves.
\end{remark}

\section{Homogeneity}

Theorem~\ref{thm:pointwise} alone does not yield a dilator. Girard's
output $F+F''$ must act on embeddings, and the difficulty is visible in the
kind III case: the order $V_1$ at which the negative premise is read
depends on the output $E_1$. To compare two evaluations we must therefore
compare the outputs of the positive premise first, and then use the
resulting map as part of the parameter map of the negative premise.

\paragraph{Related inputs.} Two input families $o,o'$ are \emph{related}
along $c:\Cst\to\Cst'$ and $\tau_Z:kZ\to k'Z$ (\tt{RelIn}) if $\tau$
preserves top, the end sequents and the rules, and
$\tau_Z\sigma^{-1}(\mathrm{inl}\,w)=\sigma'^{-1}(\mathrm{inl}((c+\id)w))$.
For an embedding $f:Y\to Y'$ write $\beta_f=\tau_{Y'}\circ k(f)$.

\paragraph{Relation of derivations.} For $u:A+E\to A'+E'$, a cut-free
derivation $d'$ \emph{restricts} to $d$ along $u$ (\tt{RelF}) if $d'$ is
literally the image of $d$: the same rule at each node, with labels,
points and environments mapped by $u$.

\begin{theorem}[\tt{BetaIDS4HomFin.Hom.hom}]\label{thm:hom}
Let $o,o'$ be related along $(c,\tau)$, with certificates $ac,ac'$. Let
$f:Y\to Y'$ be an embedding of well-orders, and let $(E,d)$ and $(E',d')$
be the outputs of Theorem~\ref{thm:pointwise} at $(o,ac,Y)$ and
$(o',ac',Y')$. There is an explicitly computed embedding $e:E\to E'$ such
that $d'$ restricts to $d$ along $\beta_f+e$.
\end{theorem}

The proof is a second recursion on the certificate, with the target
certificate split in step. In each case $e$ is assembled from the maps of
the premises:
\begin{itemize}
\item kept rules: $e=\sum_a e_a$ over the label map induced by $\beta_f$
(\tt{TrCSquare.DSumMap});
\item $(\overline r)$: $e_2=e_1+\id_1$;
\item kind III: $e_1:E_1\to E_1'$ comes from the positive premise at
$(Y,Y',f)$. Then
\[
 f_1 \;=\; f+\bigl((1+t)+e_1\bigr)\;:\;V_1\to V_1',
\]
where $t$ is the upper-tail map read off $\beta_f$. The negative premise
is compared at $(V_1,V_1',f_1)$, which gives $e_2$, and
$e_3=e_1+((1+t_1)+e_2)$.
\end{itemize}
The kind III case is where homogeneity is genuinely needed. The parameter
map $f_1$ of the second comparison is built from the output map $e_1$ of
the first, so the theorem must be proved for arbitrary embeddings and
arbitrary related inputs, not only for the identity. The cut itself
commutes with the maps (\tt{S2CRel}, \tt{BetaIDS4HomRen}), as do the
renamings 11.4.5 and 11.4.7.

\section{Certificates, identity and inverse laws}

\begin{theorem}[\tt{BetaAcc.AccR.accProp}]\label{thm:accprop}
For the generic accessibility encoding, any two certificates of the same
object are equal.
\end{theorem}

A certificate at $x$ is a tree $\mathrm{node}\,x_0\,fs$ with a proof of
$x_0=x$ and the certificates of the children. The pair consisting of $x_0$
and this proof lies in a singleton type, so the root equation is handled
by path induction, without any assumption that the type of objects is a
set. Induction on the tree identifies corresponding children, and function
extensionality identifies the two child families.

\begin{corollary}[\tt{BetaIDS4DiagIso.At.outputEq}]
For fixed $o$ and $Y$, the outputs of Theorem~\ref{thm:pointwise} at any
two certificates are equal. No further hypothesis is needed.
\end{corollary}

So once accessibility is obtained, the choice of certificate is
immaterial. The comparison maps of Theorem~\ref{thm:hom} are a separate
matter: they are computed by recursion, and equality of the outputs does
not by itself say what the maps are.

\begin{theorem}[\tt{BetaIDS4Diag.Fin.diagId}]\label{thm:diag}
Assume the pointwise identity law $k(\id_Z)(p)=p$ for $o$. Let $c$,
$\tau_Z$ and $f:Y\to Y$ be pointwise identities and let the relation
$rl:\mathrm{RelIn}(o,o,c,\tau)$ be arbitrary. Then the first component of
$\mathrm{hom}(o,o,c,\tau,rl,ac,ac,Y,Y,f)$ fixes every point of the
computed extension.
\end{theorem}

The relation $rl$ must be arbitrary because the recursive calls construct
their own relations (shifted and plain predecessors). The statement is
about the map that $\mathrm{hom}$ computes, not about a
re-implementation, and this has a cost: the recursion is \emph{not}
diagonal definitionally, even when the input is. The dispatcher compares
the root rule with its $\tau$-image along a path that is not reflexivity.
The stage labels of a kept rule move from $\nu$ to $\beta\nu$, and the
$(xl)$ children are called at $c(c_0)$ and $f(y)$ rather than at $c_0$ and
$y$. The proof therefore establishes a heterogeneous law at each node: the
computed map is the transport along a path of indices, or equivalently a
path in a total space. It uses:
\begin{itemize}
\item uniqueness of identity proofs for rules (by decidable equality and
Hedberg's theorem), for Booleans, for the three-valued comparisons, and
for carriers of linear orders;
\item based path induction at an endpoint, to contract a stage
$\nu=\sigma^{-1}(\mathrm{inl}\,w)$ to the canonical one;
\item transport over a proposition, normalised to fix every point, to pass
between the summands of the kept rule and the total space of a stage;
\item for kind III, the fact that the tail maps $t,t_1$ of a pointwise
identity are identities (from injectivity of $\sigma^{-1}$, without
computing the decision that defines them).
\end{itemize}
Function extensionality enters only through equality of maps of linear
orders. The km identity law is a clean sufficient hypothesis. Whether it
is necessary is open: on invalid branches and on stage labels above top
the construction puts the empty order, so necessity is not established on
every label.

\begin{corollary}[\tt{BetaIDS4DiagIso.At.Laws}]
Under the km identity law, for certificates $a,b$ of $o$, let
$e_{ab}:E_a\to E_b$ be the map computed by $\mathrm{homSame}$ at $\id_Y$.
Then $e_{ba}\circ e_{ab}=\id$, and $(e_{ab},e_{ba})$ is an isomorphism.
\end{corollary}

Transport along the equality $a=b$ of Theorem~\ref{thm:accprop} reduces
$e_{ba}(e_{ab}(z))=z$ to $e_{aa}(e_{aa}(z))=z$, which is
Theorem~\ref{thm:diag} applied twice.

\section{What remains}

\paragraph{Composition.} For embeddings $f:Y\to Y'$ and $g:Y'\to Y''$ and
three certificates, does the map computed for $g\circ f$ agree with the
composite of the two computed maps? This is the next step towards
Girard's functorial output dilator $F+F''$, together with the trace-range
laws of 11.4.4(iii). We expect the same heterogeneous technique to apply,
with the additional difficulty that in the kind III case the composite
parameter map $f_1$ must itself be compared with a composite. This is not
done.

\paragraph{Kind II cuts.} Girard removes kind II cuts first (11.4.3) and
then applies 11.4.4. For families, the output of that elimination would
have to be an input family again, with an accessibility certificate built
by the same recursion, since its subtrees are not subtrees of the input.
For cuts \emph{below top} this detour is not needed, and the theorems
above already include them. Such a cut is kept as a kind II cut at
$(kY+E,\mathrm{inl}\,\mathrm{top}_Y)$ by $\mathrm{rep}$
(\texttt{RepR.stLow}), and the kept-rule node ends with the pointwise
elimination of kind II cuts (\texttt{BetaIDCfCut}) on the assembled rule;
the premises are input predecessors, so the induction on accessibility is
unaffected. In the homogeneity proof the two sides cannot disagree on the
kind of a cut, since the target environment is the image of the source
one (\texttt{NodeD.lowBar}, \texttt{barLow}).

Atomic cuts at stages $\geq\mathrm{top}$ are covered as well. $\mathrm{rep}$
turns such an occurrence into $\Phib$, and the kind III construction needs
only that the \emph{represented} occurrence be $\Phib$
(\texttt{K3Node} takes \texttt{RepBar}). For homogeneity, an embedding
preserving top preserves and reflects this test on corresponding atoms
(\texttt{NodeD.rbPres}, \texttt{rbRefl}), so both sides classify a cut
alike. The boundary case is checked end to end
(\texttt{BetaIDS4Bound}): a concrete accessible family whose root is a cut
on $I^{\mathrm{top}_Y}(t)$, both premises closing on a true arithmetic
atom, satisfies the input condition but not the below-top test, and the
pointwise theorem, homogeneity and the diagonal law apply to it.

What remains are \emph{compound cuts with at least one occurrence
represented as $\Phib$}: mixed formulas such as $I^\mu(t)\wedge I^\nu(u)$
with $\mu<\mathrm{top}\le\nu$, represented as $I^\mu(t)\wedge\Phib(u)$, and
also formulas such as $\Phib(t)\wedge\Phib(u)$. They need a compound-cut
argument. Until it is given, the strength of unrestricted elimination is
open.

\paragraph{Where the strength sits.} Under the restrictions above,
everything after the certificate is constructive: normalisation,
homogeneity, certificate independence and the identity and inverse
laws. The certificate itself comes from correctness by a bar argument with
excluded middle and dependent choice. This is consistent with Girard's
11.6.4, by which full cut elimination implies $\Pi^1_1$-$\mathrm{CA}$. The
cut-free rules introduce only subformulas when read from conclusion to
active premise, allowing number instances and the prescribed unfoldings
$\Phi(I^\nu,t)$ and their negations in the inductive rules.

\paragraph{Status.} All theorems stated here are checked in Unit; the
development has no holes and no postulates. The two standing hypotheses,
the cut condition (no compound cut with an occurrence represented as
$\Phib$), and (for
Theorem~\ref{thm:diag} and its corollary) the km identity law are as
stated. The remark on composition is an expectation, not a result.

\begin{thebibliography}{9}
\bibitem{ptlc2} J.-Y.~Girard, \emph{Proof Theory and Logical Complexity},
vol.~II, chapter~11 (\S11.4: cut elimination for the inductive logic;
11.6.4: the reversal).
\bibitem{am} R.~Akiyoshi and G.~Mints, \emph{An extension of the
omega-rule}, Archive for Mathematical Logic \textbf{55} (2016),
nos.~3--4, 593--603. \href{https://doi.org/10.1007/s00153-016-0482-y}{doi:10.1007/s00153-016-0482-y}.
\end{thebibliography}

\end{document}
